% t1e_A.tex — Topic1ExplanationsPart1, pp. 1–9:
% front matter, Topic 1-1 Questions 1–4, Topic 1-2 Questions 1–3.
% Fragment: \input into a wrapper that loads mathdocs.sty.

\begin{center}
  {\LARGE\bfseries Explanations for Independent Practice (Part 1 Only)}\par
  \medskip
  {\large Topic 1 --- Written Explanation Companion}
\end{center}
\medskip

\noindent This companion restates each Part 1 question, shows the checked
answer, and explains the work very slowly. Follow the steps in order. Read
every Important words list before you start the numbered steps.

\blockhead{How to use this companion (read once)}
\noindent Instructor support, online resources, and AI are allowed on
\textbf{explanations only}. On Independent Practice Parts 1 and 2, do
\textbf{not} use the internet or AI. You may get \textbf{limited in-person
help} only. These explanations prepare you for Part 1. If you can do Part 1 on
your own, you can do Part 2. When you miss a problem, \textbf{circle} it, get
help, then \textbf{rework} until you can solve it independently and correctly.

% ---------------------------------------------------------------------------
\lesson{Topic 1-1 Relate Integers and Opposites}

\question{Question 1}

\blockhead{Question}
\noindent A climber climbs $4.5$ meters up a rock wall. Name the move that is
the additive inverse of that climb, and explain why the two moves add to zero.

\checked{$-4.5$ meters ($4.5$ meters downward); $+4.5 + (-4.5) = 0$}

\blockhead{What the question is asking}
\noindent This question asks for the move that undoes a climb of $4.5$ meters
completely (same size, opposite direction, so the two moves add to zero).

\blockhead{Important words}
\begin{words}
  \item \term{Nonzero}means not equal to zero. The symbol $\neq$ means ``is not equal to.''
  \item A \term{fraction}$\dfrac{a}{b}$ means $a$ equal parts out of $b$ equal parts ($b \neq 0$).
  \item \term{Tenths}are equal parts made when one whole is split into $10$ parts.
        \term{Hundredths}are equal parts made when one whole is split into $100$ parts.
  \item A \term{decimal}uses a decimal point to show parts of a whole (tenths, hundredths, $\ldots$).
  \item A \term{whole number}is $0, 1, 2, 3, \ldots$. It has no fraction or decimal part.
  \item A \term{terminating}decimal stops, like $0.5$ or $4.5$.
  \item A \term{repeating}decimal goes forever in a pattern, like $0.333\ldots$.
  \item An \term{integer}is one of the numbers $\ldots, -2, -1, 0, 1, 2, \ldots$; it has no decimal part.
  \item A \term{rational number}can be written as a fraction $\dfrac{a}{b}$ (with $b \neq 0$) or as a
        terminating or repeating decimal. The decimal $4.5$ is rational because $4.5 = \dfrac{9}{2}$.
  \item A \term{positive number}is greater than zero. A \term{negative number}is less than zero.
  \item A \term{signed}number (or signed move) uses a plus or minus sign to show direction:
        $+$ means positive (up / right / gain), and $-$ means negative (down / left / loss).
  \item To \term{cancel}(here) means two moves or numbers undo each other completely and leave zero.
  \item An \term{additive inverse}of a number is the number that has the \textbf{same size} but the
        \textbf{opposite direction}, so the two numbers add to zero (they cancel).
  \item \term{Vertical}means up and down.
  \item A \term{climb}of $4.5$ meters means moving \textbf{up} $4.5$ meters. We can write that as
        $+4.5$ meters.
  \item \term{Downward}means the opposite direction of up.
  \item To \term{flip}the direction means change up$\leftrightarrow$down (or $+ \leftrightarrow -$)
        while keeping the same size.
\end{words}

\blockhead{Work the problem one tiny step at a time}
\begin{steps}
  \item Write the climb as a signed move: $+4.5$ meters (up). The $+$ sign shows the positive /
        upward direction.
  \item The additive inverse must use the \textbf{same distance}: $4.5$ meters.
  \item The additive inverse must use the \textbf{opposite direction}: down (negative).
  \item So the inverse move is $-4.5$ meters, which means $4.5$ meters downward.
  \item Check: $+4.5 + (-4.5) = 0$. The two moves cancel. That is what additive inverse means.
\end{steps}

\finalans{$4.5$ meters downward}

\blockhead{Why this makes sense}
\noindent The climb and its inverse cancel because they have the same size and opposite signs.

% ---------------------------------------------------------------------------
\question{Question 2}

\blockhead{Question}
\noindent On a number line, zero is located in the middle. We have two mystery
numbers, $a$ and $b$. We are told that $a < 0$, $b > 0$, and $a + b = 0$. Write
an equation about the absolute values of $a$ and $b$ that must be true.

\checked{$\abs{a} = \abs{b}$}

\blockhead{What the question is asking}
\noindent This question asks for a math sentence with an equals sign that
compares how far $a$ and $b$ are from zero (distance from zero, which is never
negative).

\blockhead{Important words}
\begin{words}
  \item \term{Negative}means less than zero (left of zero). The symbol $<$ means ``less than.''
  \item \term{Positive}means greater than zero (right of zero). The symbol $>$ means ``greater than.''
  \item A \term{number line}is a straight line with numbers in order. Zero is in the middle.
        Negative numbers are left of zero. Positive numbers are right of zero.
  \item A \term{variable}is a letter that stands for a number. Here $a$ and $b$ are variables.
  \item An \term{equation}is a math sentence with an equals sign that says two sides are equal.
  \item An \term{additive inverse}of a number is the number with the same size but the opposite
        direction, so the two numbers add to zero.
  \item To \term{cancel}(here) means two numbers undo each other completely and leave zero
        (they are additive inverses).
  \item \term{Distance}tells how far apart two points are. Distance is never negative.
  \item \term{Absolute value}of a number, written \barsym, is the distance from that number to zero.
        Distance is never negative.
\end{words}

\blockhead{Work the problem one tiny step at a time}
\begin{steps}
  \item We are told $a + b = 0$.
  \item That means $b$ is the additive inverse of $a$: $b = -a$.
  \item So $a$ and $b$ sit on \textbf{opposite sides} of zero.
  \item Because they cancel (add to zero), they must be the \textbf{same distance} from zero.
  \item Distance from zero is absolute value. So $\abs{a}$ equals $\abs{b}$.
  \item Write the equation: $\abs{a} = \abs{b}$.
  \item Numeric check: if $a = -5$ and $b = 5$, then $a + b = 0$, and $\abs{-5} = 5 = \abs{5}$.
\end{steps}

\finalans{$\abs{a} = \abs{b}$}

\blockhead{Why this makes sense}
\noindent If two numbers add to zero, they are additive inverses (opposites).
Opposites are always the same distance from zero, so their absolute values
match. A common mistake is writing $a = b$ (false, because one is negative and
one is positive). Check: $-5 + 5 = 0$ and $\abs{-5} = \abs{5}$.

% ---------------------------------------------------------------------------
\question{Question 3}

\blockhead{Question}
\noindent Evaluate: $\abs{-9} = $

\checked{$9$}

\blockhead{What the question is asking}
\noindent This question asks how far $-9$ is from zero on a straight line of
numbers (distance from zero is never negative).

\blockhead{Important words}
\begin{words}
  \item An \term{operation}is a math action. Add joins amounts. Subtract takes away or finds a
        difference. Multiply makes equal groups. Divide splits into equal groups.
  \item An \term{expression}is a math phrase made of numbers, operations, and sometimes symbols
        like absolute-value bars (it does not have to be a full equation).
  \item A \term{decimal}uses a decimal point to show parts of a whole.
  \item To \term{evaluate}means to find the value of an expression.
  \item A \term{negative number}is less than zero. A \term{positive number}is greater than zero.
  \item A \term{number line}is a straight line with numbers in order. Zero is in the middle.
        Negatives are left; positives are right.
  \item \term{Distance}tells how far apart two points are. Distance is never negative.
  \item \term{Absolute value}\barsym\ means distance from zero on the number line. Distance is
        never negative.
  \item An \term{integer}is one of the numbers $\ldots, -2, -1, 0, 1, 2, \ldots$; it has no decimal part.
  \item A \term{unit}on the number line is one tick-mark step from one integer to the next.
  \item The number $-9$ is nine units left of zero.
  \item The absolute-value bars do \textbf{not} mean ``make the number negative.'' They report how
        far the number sits from zero.
\end{words}

\blockhead{Work the problem one tiny step at a time}
\begin{steps}
  \item Picture a straight line of numbers. Mark zero in the middle.
  \item Find $-9$. It is $9$ steps to the left of zero.
  \item Count the steps from $0$ to $-9$ one by one: $1, 2, 3, 4, 5, 6, 7, 8, 9$. That is $9$ steps.
  \item Distance does not care about left or right. Keep the count only: $9$.
  \item Absolute value writes that distance with bars: $\abs{-9} = 9$.
  \item Check another way: absolute value of a negative number is the positive version of that
        number's size, so $\abs{-9}$ matches the size $9$.
  \item Contrast check: $\abs{9} = 9$ also (same distance on the right side). Distance from zero is
        $9$ either way.
\end{steps}

\numline{-9}{0}{3}{-9}

\finalans{$9$}

\blockhead{Why this makes sense}
\noindent Absolute value never gives a negative answer. The bars turn $-9$ into
the positive distance $9$. Check: nine steps left of zero is still a distance of
$9$. A common mistake is thinking $\abs{-9} = -9$.

% ---------------------------------------------------------------------------
\question{Question 4}

\blockhead{Question}
\noindent Order from least to greatest: $-3$, \ $8$, \ $-1$

\checked{$-3$, \ $-1$, \ $8$}

\blockhead{What the question is asking}
\noindent This question asks you to list the three numbers from smallest to
largest.

\blockhead{Important words}
\begin{words}
  \item \term{Least to greatest}means smallest first, then bigger, then biggest.
  \item On a \textit{number line}, numbers farther to the \textbf{left} are smaller. Numbers farther
        to the \textbf{right} are larger.
  \item A \term{negative}number is less than zero. A \term{positive}number is greater than zero.
  \item \term{Absolute value}means distance from zero. Ordering by absolute value alone is wrong
        because it ignores which side of zero a number is on.
\end{words}

\blockhead{Work the problem one tiny step at a time}
\begin{steps}
  \item Draw a number line and mark $-3$, $-1$, and $8$.
  \item Look which mark is farthest left: $-3$. That is the least.
  \item Next from the left: $-1$.
  \item Farthest right: $8$. That is the greatest.
  \item Write them in that order: $-3$, \ $-1$, \ $8$.
\end{steps}

\numline{-4}{9}{1}{-3,-1,8}

\finalans{$-3$, \ $-1$, \ $8$}

\blockhead{Why this makes sense}
\noindent Left means smaller. $-3$ is left of $-1$, and both negatives are left
of positive $8$. A common mistake is ordering by absolute value and writing
$-1$, \ $-3$, \ $8$.

% ---------------------------------------------------------------------------
\lesson{Topic 1-2 Understand Rational Numbers}

\question{Question 1}

\blockhead{Question}
\noindent In a survey, $3$ out of $8$ students chose basketball as their
favorite sport. What is the decimal equivalent of the fraction of students who
chose basketball?

\checked{$0.375$}

\blockhead{What the question is asking}
\noindent This question asks you to change $\dfrac{3}{8}$ into a number that
uses a dot for parts of a whole.

\blockhead{Important words}
\begin{words}
  \item A \term{fraction}$\dfrac{3}{8}$ means $3$ parts out of $8$ equal parts. The top number is
        the \textit{numerator}. The bottom number is the \textit{denominator}.
  \item A \term{decimal}uses a decimal point to show parts of a whole, like $0.375$.
  \item \term{Equivalent}means ``equal in value.''
  \item \term{Long division}is a written way to divide one digit at a time and keep track of what
        is left.
  \item A fraction bar means \textit{divide}: top $\div$ bottom.
  \item A \term{remainder}is what is left over after a division step.
  \item To \term{bring down}a digit means move the next digit (often a $0$) into the current
        remainder so you can keep dividing.
  \item A \term{decimal digit}is a digit to the right of the decimal point (after the dot).
\end{words}

\blockhead{Work the problem one tiny step at a time}
\begin{steps}
  \item Write the fraction: $\dfrac{3}{8}$.
  \item Divide: $3 \div 8$. Since $3$ is smaller than $8$, write $3.000$.
  \item $8$ goes into $30$ three times: $8 \times 3 = 24$. Subtract: $30 - 24 = 6$. Remainder $6$.
  \item Bring down $0$ to make $60$. $8$ goes into $60$ seven times: $8 \times 7 = 56$.
        Subtract: $60 - 56 = 4$. Remainder $4$.
  \item Bring down $0$ to make $40$. $8$ goes into $40$ five times: $8 \times 5 = 40$.
        Subtract: $40 - 40 = 0$. Remainder $0$.
  \item The decimal digits are $3$, then $7$, then $5$. So $3 \div 8 = 0.375$.
\end{steps}

\finalans{$0.375$}

\blockhead{Why this makes sense}
\noindent Long division stops with remainder $0$, so the decimal terminates
(stops; no leftover repeating pattern) at $0.375$. Check: $0.375 \times 8 = 3$.

% ---------------------------------------------------------------------------
\question{Question 2}

\blockhead{Question}
\noindent A student uses long division to convert a fraction $\dfrac{x}{y}$,
where $y \neq 0$, into a decimal. Which outcome is \textbf{NOT} possible?

\begin{enumerate}[label=\Alph*), leftmargin=2.6em, itemsep=0.35\baselineskip, topsep=0.4\baselineskip]
  \item The decimal terminates (stops).
  \item The decimal repeats a pattern.
  \item The decimal neither terminates nor repeats.
\end{enumerate}

\checked{C}

\blockhead{What the question is asking}
\noindent This question asks which outcome cannot happen when you divide the
numerator of a fraction by its nonzero denominator.

\blockhead{Important words}
\begin{words}
  \item \term{Long division}is a written way to divide one digit at a time and keep track of what
        is left.
  \item A \term{fraction}$\tfrac{x}{y}$ means the top number $x$ is divided by the bottom number $y$.
  \item The \term{numerator}is the top number being divided. Here it is $x$.
  \item The \term{denominator}is the bottom number that divides the top number. Here it is $y$.
  \item \term{Nonzero}means not equal to zero. The denominator must be nonzero because division by
        zero is not allowed.
  \item A \term{whole number}is one of the numbers $0, 1, 2, 3, \ldots$.
  \item A \term{decimal}uses a decimal point to show parts smaller than one whole.
  \item A \term{terminating decimal}stops, like $0.5$ or $0.375$.
  \item A \term{repeating decimal}goes on forever but in a repeating pattern, like $0.333\ldots$.
  \item A \term{quotient}is the answer to a division problem.
\end{words}

\blockhead{Work the problem one tiny step at a time}
\begin{steps}
  \item When you divide two whole numbers (nonzero denominator), the decimal \textbf{must} either
        terminate or repeat.
  \item Statement A can be true for some fractions (example: $\tfrac{1}{2} = 0.5$).
  \item Statement B can be true for some fractions (example: $\tfrac{1}{3} = 0.333\ldots$).
  \item Statement C says the decimal neither terminates nor repeats. That outcome is
        \textbf{not possible} for a fraction of whole numbers.
  \item So the outcome that is NOT possible is C.
\end{steps}

\finalans{C}

\blockhead{Why this makes sense}
\noindent Fractions of whole numbers always become terminating or repeating
decimals. Statement C describes something that does not happen for those
fractions.

% ---------------------------------------------------------------------------
\question{Question 3}

\blockhead{Question}
\noindent In a class debate, Jordan says $0.\overline{27}$ is rational because
its digits repeat. Casey says it is not rational because it never ends. Which
student is correct? Use a fraction or equation to support the claim.

\checked{Jordan; $0.\overline{27} = \dfrac{3}{11}$.}

\blockhead{What the question is asking}
\noindent Decide whether a repeating decimal is rational and prove it.

\blockhead{Important words}
\begin{words}
  \item A repeating decimal has a digit pattern that repeats forever.
  \item A rational number can be written as a fraction of integers with a nonzero denominator.
\end{words}

\blockhead{Work the problem one tiny step at a time}
\begin{steps}
  \item Let $x = 0.\overline{27}$.
  \item Multiply by $100$: $100x = 27.\overline{27}$.
  \item Subtract: $99x = 27$, so $x = \dfrac{27}{99} = \dfrac{3}{11}$.
\end{steps}

\finalans{Jordan is correct; $0.\overline{27} = \dfrac{3}{11}$.}

\blockhead{Why this makes sense}
\noindent The decimal has an exact fraction name, so it is rational.
